% ============================================================
% 附件② — 与 A11 系统关系及功能互补性说明
% 评审会六项附件待办第 2 项 | 2026-07-13 评审会
% 编译: xelatex -interaction=nonstopmode 附件2-与A11系统关系及功能互补性说明.tex
% ============================================================
\documentclass[11pt,a4paper]{ctexart}
\usepackage{xeCJK}
\usepackage{graphicx}
\usepackage{float}
\usepackage{fancyhdr}
\usepackage{hyperref}
\usepackage{geometry}
\usepackage{longtable}
\usepackage{booktabs}
\usepackage{enumitem}
\usepackage{caption}
\usepackage{array}

% 字体
\setCJKmainfont{Microsoft YaHei}

% 页面
\geometry{a4paper, margin=2.5cm}

% 页眉页脚
\pagestyle{fancy}
\fancyhead[L]{时序数据采集与应用管理系统 — 附件②}
\fancyhead[R]{与 A11 系统关系及功能互补性说明}
\fancyfoot[C]{\thepage}

\hypersetup{
  colorlinks=true, linkcolor=black, urlcolor=blue,
  pdftitle={与 A11 系统关系及功能互补性说明},
}

\begin{document}

\begin{center}
{\Huge\bfseries 与 A11 系统关系及功能互补性说明}\\[14pt]
{\large 时序数据采集与应用管理系统 — 附件②}\\[8pt]
{\large 版本 V1.0 \quad 编制日期 2026-08}
\end{center}

\vspace{1em}
\noindent\hrulefill
\vspace{1em}

% ═══════════════════════════════════════════
\section{关系定位：并行互补，非改造、非替代}

本项目（时序数据采集与应用管理系统）与既有 A11 油气生产物联网系统之间的关系可概括为：\textbf{在不改造、不影响 A11 正常运行的前提下，以路由监听方式实现时序数据双路采集，与 A11 并行运行、功能互补}。

\begin{enumerate}[leftmargin=*]
  \item \textbf{不改 A11}：本项目不修改 A11 系统任何配置、\textbf{不在 A11 核心生产服务器（PHD 实时库等）安装任何软件}、不中断 A11 既有采集链路——A11 继续承载原有生产业务，本项目在其旁路获取数据副本；本系统仅以\textbf{轻量只读采集组件}形式与 IO 数据接入服务器共存部署（见 4.1 利旧部署），不影响 A11 服务。
  \item \textbf{保留现有全部配置与业务系统集成}：A11 既有配置（已建点位、量程频次、物模型、A2/A5 资产体系）与既有业务系统的集成关系（数据交换、接口对接）全部原样保留、不中断、不重建——A11 侧零配置变更、业务系统零感知，本项目以主数据映射方式对接既有点位（见第 5 节），以数据湖门户嵌入方式融入业务系统生态（见附件⑥）。
  \item \textbf{不替代 A11}：A11 是集团统建系统，承担油田生产数据采集的既有职能；本项目定位为数据中台实时数据采集子模块，与数据湖逻辑集成，二者各自职责清晰、互为补充。
  \item \textbf{两套体系合一}：本项目采集的时序数据全部统一注册至数据湖主数据体系，点位管理、物模型构建以主数据为唯一源头，不形成独立于数据湖的第二套数据管理体系（评审会硬性要求之二）。
\end{enumerate}

% ═══════════════════════════════════════════
\section{A11 系统现状与能力边界}

\subsection{采集模式现状}

A11 系统基于 PHD 实时数据库两级架构运行：前端采集节点按固定周期扫描取数，经抽吸/轮询方式上送，全系统统一频率、统一配置。从油田公开口径及我方对现场 IO 服务器采集规律的实测（生产网只读观测，未安装任何软件）看：

\begin{itemize}[leftmargin=*]
  \item 数据为\textbf{分钟级定时抽吸}（约 10--30 分钟一档），功图类曲线数据约 30 分钟/副；
  \item 报警依赖后台批量处理，无本地流式计算，无法做到秒级/亚秒级异常响应；
  \item 协议支撑有限，视频流、智能仪表、新能源等新兴设备源无法接入；
  \item 底层依赖老旧 Windows 环境，无国产 CPU/OS 适配能力，与信创替代要求存在差距。
\end{itemize}

\subsection{能力边界对生产需求的制约}

油田近年推进电参数字化、工况智能诊断、设备健康预警等应用，其算法能力在局地试点已获验证，但受限于分钟级采集底座：高频异常（皮带断裂、电潜泵参数突变、频率跃升冲击）在数据刷新间隙内无法被识别；差异化报警（逐井、逐工况）无逐井调优通道；动态接入（移动网关、设备增减）无统一感知机制。\textbf{瓶颈不在算法，而在数据密度与采集灵活性}——这正是本项目互补价值的切入点。

% ═══════════════════════════════════════════
\section{功能互补矩阵}

\begin{table}[H]\centering
\caption{与 A11 的功能互补对照（评审会认可的必要性依据细化）}
\begin{tabular}{p{2.6cm} p{4.2cm} p{4.2cm} p{2.6cm}}
\toprule
\textbf{能力维度} & \textbf{A11 现状} & \textbf{本项目补充能力} & \textbf{对应模块} \\
\midrule
采集频率 & 分钟级定时抽吸 & 可配置 500ms--60s，异常自动加密（10min→秒级）、低峰降频 & 1.2 定制动态采集 \\
边缘计算 & 无本地流式引擎，报警依赖后台批处理 & 入库前流水计算，15 种内置算法 + 每作业区定制算法，单条 <1ms & 1.4 流式计算引擎 \\
协议兼容 & 有限工业协议 & MQTT/HTTP/UDP/RTSP 多协议，10 家网关厂家适配，非标协议按需接入 & 1.2 定制动态采集 \\
信创适配 & 依赖 Windows 老旧环境 & 全栈开源组件 + 自主定制，麒麟/飞腾等信创环境适配 & 一-9 全链路安全与审计追溯平台 \\
跨网贯通 & 生产网/办公网隔离 & 生产网边缘采集 + 办公网 DMZ 短期缓存入湖，安全通道穿透 & 1.5 实施交付 \\
\bottomrule
\end{tabular}
\end{table}

\textbf{一句话概括互补关系}：A11 管住存量（统建标准采集、既有生产业务不变），本项目管住增量（高频、动态、差异化、信创）——两条链路并行，数据汇聚于数据湖主数据体系。

% ═══════════════════════════════════════════
\section{技术实现：路由监听双路采集}

\subsection{部署形态（评审会硬性要求之一）}

在 IO 服务器侧以\textbf{路由监听}方式获取时序数据，实现双路采集：\textbf{原 A11 采集链路保持原样运行，本项目在旁路监听同一数据流}，实时同步。

\begin{itemize}[leftmargin=*]
  \item \textbf{采集形态按终端接入类型分流}：
  \begin{itemize}
    \item \textbf{静态 IP 终端（有线/固定链路，设备侧为服务端）}：采用\textbf{旁路方案}——数据端点在设备侧（server），本系统以\textbf{客户端身份主动连接采集}：不改设备任何配置、设备侧无感知（服务端仅多一个客户端连接）、原 A11 采集链路原样不动；
    \item \textbf{动态 IP 终端（4G/移动网关，设备侧为客户端）}：只能采用\textbf{桥接方案}——设备主动建连至原服务器（client$\rightarrow$server），不改设备层 IP，\textbf{则必须将原服务器 IP 与端口抢占做桥接}：桥接层绑定原 IP:端口，设备连接行为不变，数据到达后复制一份并透明转发原 A11 服务——唯此一路径，方能对动态终端全覆盖；
    \item 两种形态\textbf{按终端类型并存}：静态 IP 部分走客户端主动采集、动态 IP 部分走桥接转发，统一汇聚入采集链路；实施前以现场调研核实静态/动态终端占比，两种形态按现场条件组合实施。
  \end{itemize}
  \item \textbf{硬件投入与利旧部署}：项目新增硬件仅 2--3 台服务器（生产网边缘计算 1 台、DMZ 短期缓存 1 台、办公网汇聚中枢 1 台，可按现场合并为 2 台）；其余设施（IO 服务器、网络、存储、既有设备）\textbf{全部利旧}——本系统采集组件\textbf{直接部署于既有 IO 服务器}（轻量组件：静态 IP 客户端采集进程 + 动态 IP 桥接转发进程，CPU 占用 $<$5\%、延迟增加 $<$20$\mu$s），与 A11 采集服务同机共存：独立进程、独立端口、资源限额、故障自动静默，不改变 A11 任何行为；
  \item \textbf{接入节奏}：三阶段渐进——观察（bypass，仅只读观察、不介入）$\rightarrow$ 桥接复制（mirror，桥接层绑定原地址端口、复制并原样转发，原链路业务不断）$\rightarrow$ 稳定并行（stable），任一阶段异常可即时回退，原链路零影响；\textbf{抢占端口不等于接管数据——受"不改 DTU 配置、不影响 A11 原有业务"约束，A11 设备侧数据流只能以桥接模式获取：桥接层占用原地址端口仅作透明过手，\textbf{数据流始终全量流向 A11}，本项目仅复制一份副本}；
  \item \textbf{对生产系统零影响}：并联而非串联（无阻断点）、不改 A11 配置、不注入任何数据（桥接仅原样续传原链路数据，不产生、不修改任何业务内容）、故障不影响主链路。
\end{itemize}

\subsubsection{桥接模式的技术难度（高难度要求）}

在"不改 DTU 配置、不改变任何端点收发行为"的硬约束下，动态 IP 终端只能以\textbf{抢占原地址端口做透明桥接}方式获取数据，这是本项目核心技术难点之一：

\begin{itemize}[leftmargin=*]
  \item \textbf{原 IP 与端口抢占（核心难点）}：桥接层\textbf{接管原服务器 IP 与端口}——原 A11 服务迁至备用地址/端口，切换需与甲方协调、选业务低峰窗口执行，全程可回退；
  \item \textbf{透明应答与转发}：桥接层代替原服务完成 TCP 建连应答，数据复制一份后\textbf{原样转发}原 A11 服务，设备连接行为零感知——不产生、不修改任何业务内容；
  \item \textbf{TCP 流完整重组}：镜像流量需按会话重组为完整应用层数据（乱序、重传、分段、粘连），才能还原 A11 应用层帧，重组错误即数据错位；
  \item \textbf{时序与顺序严格保持}：副本与源流一致，供后续流式计算与历史回放使用，顺序错乱将导致计算失真；
  \item \textbf{长连接会话跟踪}：A11 链路为长连接心跳会话（设备重连、端口漂移、会话重建），需持续跟踪识别；\textbf{动态 IP 终端（4G 移动网关）只能以桥接方式采集——桥接是本项目主形态而非边角场景}（静态 IP 终端走客户端主动采集）；
  \item \textbf{故障完全隔离与快速回退}：桥接层任何故障不影响业务数据续传——原 A11 服务秒级切回原 IP:端口，链路自动恢复；不占资源、异常自动静默。
\end{itemize}

对应工作量已在附件④ 显式单列：路由监听分析（A11 桥接链路）45 人天（模块二）+ A11 专有协议帧解析 45 人天（模块一）+ 试点联调 A11 桥接验证 50 人天（模块八）。

\subsubsection{桥接切换操作纪律}

桥接层启用（抢占原地址端口）遵循\textbf{三条操作纪律}，作为实施与验收的共同依据：

\begin{itemize}[leftmargin=*]
  \item \textbf{切换窗口期要非常短}：切换仅在业务低峰窗口执行，窗口压缩至秒级（目标毫秒级），期间原业务流可感知的中断趋近于零；
  \item \textbf{切换前业务流摸清}：切换前对原有业务流\textbf{全面梳理建档}——设备清单、连接会话、端口占用、心跳保活机制、流量特征逐项核实，切换方案按业务流清单逐项核对，不留盲区；
  \item \textbf{异常即时回退}：切换全程值守，任何异常按预案\textbf{即时回退}——原服务秒级切回原 IP:端口、链路自动恢复，并出具切换前后业务流对比记录验证无损。
\end{itemize}

\subsection{对原有业务的零影响保障与验收口径}

本项目承诺：\textbf{实施全周期 A11 原有业务零中断、零配置变更、零性能劣化、零数据污染}，A11 继续保留运行、继续承载原有生产业务，本项目仅在其旁路并行获取数据副本。

\begin{itemize}[leftmargin=*]
  \item \textbf{机制保障}：并联而非串联（无阻断点）、不修改 A11 任何配置、不向 A11 注入任何数据（桥接仅原样续传、不产生不修改任何业务内容）、故障不影响主链路；三阶段（观察 $\rightarrow$ 桥接复制 $\rightarrow$ 稳定并行）任一阶段异常可即时回退，原链路零影响；
  \item \textbf{监控保障}：联调期间对 A11 采集链路运行状态进行并行观察记录，与本项目实施前基线比对，任何偏差即时回退并核查；
  \item \textbf{验收口径}：以"实施前后 A11 运行对比记录"作为验收项——A11 业务零中断、零配置变更、性能无劣化（采集成功率、链路时延与实施前基线一致）；桥接切换的窗口时长、切换前业务流摸底清单、回退演练记录一并随 M4 联调报告提交。
\end{itemize}

\subsection{合规性论证}

以只读方式获取数据在法规与标准层面有明确依据：

\begin{enumerate}[leftmargin=*]
  \item \textbf{合法性基础}：《网络安全法》第二十一条要求网络运营者"监测、记录网络运行状态"——采集生产数据以支撑安全生产监测，属法定职责而非越权行为；
  \item \textbf{国标认可的部署形态}：GB/T 20945--2023《网络安全审计产品技术规范》附录 A 将审计产品"并联部署（网络数据镜像或分流）"与串联部署并列，为标准认可的旁路取数据正规形态；
  \item \textbf{只读获取是行业标准做法}：NIST SP 800--82r3《OT 安全指南》（2023）明确"OT 环境中基于网络的监测能力通常经 SPAN 端口或被动网络 TAP 部署"，并强调以不负面影响运行性能与安全为前提；IEC 62443-3-3 SR 6.1 亦要求审计日志经授权只读访问——本项目获取通道只读、不写、不注入，与之一致；
  \item \textbf{等保义务}：GB/T 22239--2019（等保 2.0）要求在网络边界与重要节点实施安全审计，"审计措施缺失"属高风险项——本项目以只读授权方式获取数据，同时也是协助履行等保审计义务；
  \item \textbf{合规边界}：甲方书面授权 + 合同范围约定 + 数据归属甲方 + 最小化留存（仅保留与物模型匹配的测点数据）加密脱敏处置。
\end{enumerate}

% ═══════════════════════════════════════════
\section{数据体系：主数据唯一源头}

\begin{enumerate}[leftmargin=*]
  \item \textbf{物模型注册}：点位管理、物模型构建（G1--G8 油水井物模型，64 点）全部以数据湖主数据为唯一源头，本项目从主数据拉取设备/点位定义，不自行建立第二套体系；
  \item \textbf{自动配点}：基于物模型地址对齐自动生成配点表，减少人工维护，保证与主数据一致性；
  \item \textbf{数据流向}：采集数据 $\rightarrow$ 边缘流式计算 $\rightarrow$ DMZ 短期缓存 $\rightarrow$ 数据湖统一长期存储——本项目是数据湖实时数据采集子模块，管理界面嵌入数据湖门户，不新建孤岛。
\end{enumerate}

% ═══════════════════════════════════════════
\section{两层部署架构（评审会硬性要求之三）}

\begin{itemize}[leftmargin=*]
  \item \textbf{生产网侧}（IO 服务器旁路）：仅承担边缘采集与流式计算，毫秒级响应，不落地长周期数据；
  \item \textbf{办公网侧}（DMZ 区）：短期缓存 + 安全接入数据湖，长期存储由数据湖统一承载；
  \item \textbf{跨网安全}：双向认证通道（国密加密），生产网只出不进，控制指令穿透有独立审批。
\end{itemize}

\vspace{1em}
\noindent\hrulefill

\begin{center}
\small 本说明随项目申报材料提交，为评审会附件待办之第 2 项。\\
编制单位：迪格(杭州)物联科技有限公司 \quad 2026-08
\end{center}

\end{document}
